\section{Lattice tools and robust subsets}\label{sec:lattice-tools}

The lattice estimates in this section serve two purposes. They count
short vectors after removing directions of small determinant, and they
show that a robust collection of cube points has a difference span of
small height. We first fix the determinant conventions needed for both
uses.

\subsection{Determinants, quotients, and rational height}

A lattice \(L\) is a discrete subgroup of a finite-dimensional
Euclidean space, considered inside its real span \(E=\Span L\).
Its rank is \(\dim E\). If \(b_1,\ldots,b_r\) is a basis of \(L\), its
\emph{determinant} is the Euclidean covolume
\[
 \det L=\det\bigl(\langle b_i,b_j\rangle_{i,j=1}^r\bigr)^{1/2}.
\]
We set \(\det\{0\}=1\). Sublattices may have smaller rank. The
saturation of a sublattice \(M\subset L\) is
\(\overline M=L\cap\Span M\); it is a lattice of the same rank, and
\[
 \det M=[\overline M:M]\det\overline M.
\]
We call \(M\) \emph{primitive}, or \emph{saturated}, when
\(M=\overline M\). A primitive sublattice has a basis that extends to
a basis of \(L\). The dual lattice is always taken in the span:
\[
 L^*=\{y\in E:\langle y,x\rangle\in\Z\text{ for all }x\in L\}.
\]
In particular, the dual of the rank-zero lattice is the rank-zero
lattice.

\begin{lemma}[Orthogonal lattice quotients]\label{lem:lattice-quotient}
Let \(L\) have rank \(r\) and span \(E\), and let \(K\subset L\) be a
primitive sublattice of rank \(k\), with \(F=\Span K\). Write
\(\pi_{F^\perp}\) for orthogonal projection from \(E\) onto
\(E\cap F^\perp\). The Euclidean realization
\[
 L/K:=\pi_{F^\perp}(L)
\]
is a lattice of rank \(r-k\), canonically isomorphic to the abstract
quotient, and
\[
 \det L=\det K\,\det(L/K),\qquad
 \det L^*=(\det L)^{-1}.
\]
The dual identities are
\[
 (L/K)^*=L^*\cap F^\perp,\qquad
 K^*=\pi_F(L^*),
\]
where \(\pi_F\) is orthogonal projection onto \(F\).
For any sublattice \(M\subset L/K\), its full preimage
\(\widetilde M=\{x\in L:\pi_{F^\perp}x\in M\}\) has rank
\(k+\rank M\) and satisfies
\[
 \det\widetilde M=\det K\,\det M.
\]
These assertions include \(k=0\), \(k=r\), and \(r=0\), with all
rank-zero determinants equal to one.
\end{lemma}

\begin{proof}
Extend a basis \(b_1,\ldots,b_k\) of \(K\) to a basis
\(b_1,\ldots,b_r\) of \(L\). The projected vectors
\(\pi_{F^\perp}b_{k+1},\ldots,\pi_{F^\perp}b_r\) form a basis of the
projected lattice. Splitting Euclidean volume into the factor in \(F\)
and the factor orthogonal to \(F\) gives the first determinant identity.
The Gram matrix of the dual basis is the inverse Gram matrix, giving
the identity for \(\det L^*\).

For \(y\in E\cap F^\perp\), the equality
\(\langle y,\pi_{F^\perp}x\rangle=\langle y,x\rangle\) gives
\((L/K)^*=L^*\cap F^\perp\). Projection of \(L^*\) to \(F\) is
contained in \(K^*\). Conversely, for \(w\in K^*\), prescribe on the
basis of \(L\) the integer values \(\langle w,b_i\rangle\) for
\(i\le k\) and zero for \(i>k\). The representing vector belongs to
\(L^*\), and its projection to \(F\) is \(w\). This proves the second
dual identity. Finally, lift a basis of \(M\) and adjoin the basis of
\(K\). The resulting basis of \(\widetilde M\) gives its rank and
determinant by the same volume factorization. Empty bases give all the
rank-zero cases.
\end{proof}

A subspace \(W\subset\R^m\) is \emph{rational} if it is spanned by
integer vectors. Its integer lattice \(\Lambda_W=W\cap\Z^m\) has
rank \(\dim W\), and its \emph{height} is
\[
 H(W)=\det\Lambda_W.
\]
Thus \(H(\{0\})=H(\R^m)=1\). We will also write
\(h(W)=\log_N H(W)\).

\begin{lemma}[Height submodularity]\label{lem:height-submodularity}
For rational subspaces \(U,V\subset\R^m\),
\[
 H(U+V)H(U\cap V)\le H(U)H(V).
\]
Consequently \(h(U+V)+h(U\cap V)\le h(U)+h(V)\).
Every rational subspace has height at least one, and, for fixed \(m\)
and \(C<\infty\), there are only finitely many rational subspaces
\(W\subset\R^m\) with \(H(W)\le C\). These statements include
rank-zero subspaces.
\end{lemma}

\begin{proof}
Put \(J=U\cap V\), and project orthogonally off \(J\).
The lattice \(\Lambda_J=\Lambda_U\cap\Lambda_V\) is primitive in
each of \(\Lambda_U,\Lambda_V,\Lambda_U+\Lambda_V\). For the last
assertion, if \(u+v\in J\) with \(u\in\Lambda_U\) and
\(v\in\Lambda_V\), then \(u,v\in J\).
The projected spans of \(U\) and \(V\) intersect trivially. If
\(\overline\Lambda_U,\overline\Lambda_V\) denote their projected
lattices, the union of bases of these lattices is therefore a basis
of their sum. Orthogonally projecting the second set of basis
vectors off the span of the first can only decrease its volume.
Hence
\[
 \det(\overline\Lambda_U+\overline\Lambda_V)
 \le \det\overline\Lambda_U\,\det\overline\Lambda_V.
\]
Lemma~\ref{lem:lattice-quotient} now gives
\[
 \det(\Lambda_U+\Lambda_V)\det\Lambda_J
 \le \det\Lambda_U\,\det\Lambda_V.
\]
The saturation of \(\Lambda_U+\Lambda_V\) is
\(\Lambda_{U+V}\), whose determinant is no larger. This proves the
height inequality, and taking \(\log_N\) proves its additive form.

For a rank-\(r\) rational subspace, wedge a basis of \(\Lambda_W\).
The resulting vector in \(\bigwedge^r\Z^m\) is determined up to sign,
has Euclidean norm \(H(W)\), and determines its supporting subspace
\(W\).
Its integer coordinates are not all zero, so its norm is at least
one. A bounded ball contains finitely many vectors of
\(\bigwedge^r\Z^m\), and there are only \(m+1\) possible ranks.
This proves the finiteness assertion. For \(r=0\), the exterior
vector is \(1\) and there is just the zero subspace.
\end{proof}

For a coordinate vector \(e_i\), both possible intersections
\(U\cap\Span(e_i)\) have height one.
Lemma~\ref{lem:height-submodularity} therefore also gives
\(H(U+\Span(e_i))\le H(U)\). The next identity expresses height
ratios as determinants of normal lattices; it will transfer height
minimization to bounds on all sublattices of a dual lattice.

\begin{lemma}[Normal lattices and height]\label{lem:normal-height}
Let \(V\subset U\subset\R^m\) be rational subspaces, and put
\(\Lambda=\Lambda_U^*\). The primitive normal lattice
\[
 \mathcal N_U(V)=\Lambda\cap V^\perp
\]
has rank \(\dim U-\dim V\), and
\[
 \pi_V(\Lambda)=\Lambda_V^*,\qquad
 \det\mathcal N_U(V)=\frac{H(V)}{H(U)}.
\]
More generally, for any sublattice \(M\subset\Lambda\), set
\(V_M=U\cap(\Span M)^\perp\). Then \(V_M\) is rational and
\[
 \det M=[\mathcal N_U(V_M):M]\frac{H(V_M)}{H(U)}
 \ge \frac1{H(U)}.
\]
For \(V=U\) the normal lattice has rank zero and determinant one;
for \(V=\{0\}\) it is \(\Lambda_U^*\). If \(M=\{0\}\), the displayed
index is one and the determinant formula gives one.
\end{lemma}

\begin{proof}
The lattice \(\Lambda_V\) is primitive in \(\Lambda_U\), so the
second dual identity in Lemma~\ref{lem:lattice-quotient} gives
\(\pi_V(\Lambda_U^*)=\Lambda_V^*\). Its kernel is
\(\mathcal N_U(V)\), a primitive sublattice of the asserted rank.
Factoring the determinant of \(\Lambda_U^*\) along this kernel gives
\[
 \frac1{H(U)}
 =\det\mathcal N_U(V)\,\det\Lambda_V^*
 =\frac{\det\mathcal N_U(V)}{H(V)}.
\]

An integer basis matrix for \(\Lambda_U\) gives a rational basis for
its dual by inverting its Gram matrix. Thus \(\Span M\), and hence
\(V_M\), is rational. The saturation of \(M\) in \(\Lambda\) is
exactly \(\mathcal N_U(V_M)\). The index formula for determinants
and the first part of the lemma give the final identity; the
inequality follows from \(H(V_M)\ge1\).
\end{proof}

\subsection{Gaussian counts and determinant minimization}

For a lattice \(L\) and \(s>0\), define its Gaussian mass by
\[
 \rho_s(L)=\sum_{x\in L}\exp\!\left(-\frac{\pi\|x\|^2}{s^2}\right).
\]
In rank zero this mass equals one. The following is the scaled form
of the Reverse Minkowski Theorem of Regev and
Stephens-Davidowitz~\cite[Theorem~1.2]{ReverseMinkowski}.

\begin{theorem}[Reverse Minkowski Theorem]\label{thm:reverse-minkowski}
Let \(L\) be a lattice of rank \(j\ge1\), and let \(T>0\). Suppose
that
\[
 \det M\ge T^{\rank M}\qquad\text{for every sublattice }M\subset L.
\]
With \(\tau_j=10(\log j+2)\), one has
\[
 \rho_{T/\tau_j}(L)\le\frac32.
\]
For a rank-zero lattice the Gaussian mass is one at every scale, so
the corresponding assertion is immediate.
\end{theorem}

\begin{proof}
Apply the cited theorem to \(T^{-1}L\) in its Euclidean span.
Every rank-\(t\) sublattice of \(T^{-1}L\) has determinant at least
one, and \(\rho_{1/\tau_j}(T^{-1}L)=\rho_{T/\tau_j}(L)\).
\end{proof}

Here \(\tau_j\) is bounded above and below by absolute multiples of
\(\log(2j)\). We will use the theorem both at its stated scale and
at larger scales, through the following consequences of Poisson
summation.

\begin{lemma}[Gaussian point counts]\label{lem:gaussian-count}
Let \(L\) be a rank-\(j\) lattice. For \(R\ge0\) and \(s>0\),
\[
 \#\{x\in L:\|x\|\le R\}
 \le \exp(\pi R^2/s^2)\rho_s(L).
\]
For \(s'\ge s>0\),
\[
 \rho_{s'}(L)\le (s'/s)^j\rho_s(L).
\]
In particular, under the hypotheses of
Theorem~\ref{thm:reverse-minkowski}, if \(s_0=T/\tau_j\), then for
every \(u\ge s_0\),
\[
 \#\{x\in L:\|x\|\le R\}
 \le\frac32\left(\frac{u}{s_0}\right)^j
       \exp(\pi R^2/u^2).
\]
When \(j=0\), the point count and every Gaussian mass are one; the
first inequality is immediate and the second is equality.
\end{lemma}

\begin{proof}
Each point of norm at most \(R\) contributes at least
\(\exp(-\pi R^2/s^2)\) to \(\rho_s(L)\), proving the first bound.
The Gaussian Poisson formula in the span of \(L\) is
\[
 \rho_s(L)=\frac{s^j}{\det L}\rho_{1/s}(L^*).
\]
The Gaussian mass is nondecreasing in its scale. Applying this fact
to \(1/s'\le1/s\) in the formula proves the second bound. The last
bound follows by using the first bound at scale \(u\), then the
second bound and Theorem~\ref{thm:reverse-minkowski}. The same
Poisson identity reads \(1=1\) in rank zero.
\end{proof}

We obtain the hypotheses of the Reverse Minkowski Theorem by removing
a sublattice that minimizes a normalized determinant. The optional
rank bound below quantifies how many directions can be removed when
all original sublattices already have a weaker determinant bound.
Extremal sublattices and canonical filtrations also occur in the
reduction theory of Euclidean lattices; see
Grayson~\cite[Section~1]{Grayson}. We need only the elementary
minimization statement below.

\begin{lemma}[Determinant-minimizing sublattice]\label{lem:det-minimizer}
Let \(L\) be a lattice and let \(T>0\). The quantity
\[
 \frac{\det M}{T^{\rank M}},\qquad M\subset L\text{ a sublattice},
\]
has a minimum, attained by a primitive sublattice \(L_0\). It
satisfies
\[
 \det L_0\le T^{\rank L_0},\qquad
 \det M\ge T^{\rank M}\quad\text{for every sublattice }M\subset L/L_0.
\]
If in addition \(T_0>T\), \(\Delta\ge0\), and
\[
 \det M\ge e^{-\Delta}T_0^{\rank M}
       \quad\text{for every sublattice }M\subset L,
\]
then
\[
 \rank L_0\le\frac{\Delta}{\log(T_0/T)}.
\]
For \(L=\{0\}\) the unique choice is \(L_0=\{0\}\), and every
assertion holds with rank zero and determinant one.
\end{lemma}

\begin{proof}
The rank-zero lattice has normalized determinant one. For a fixed
positive rank \(t\), the determinant of a sublattice is the norm of
the wedge of one of its bases in the discrete lattice
\(\bigwedge^t L\). There are only finitely many such norms at most
\(T^t\). Thus among all normalized determinants no greater than
one there are only finitely many values, and the minimum is
attained. Saturating a minimizer can only decrease its determinant
while preserving its rank, so a primitive minimizer may be chosen.
Comparison with the rank-zero lattice gives
\(\det L_0\le T^{\rank L_0}\).

For a sublattice \(M\subset L/L_0\), lift it to \(\widetilde M\)
as in Lemma~\ref{lem:lattice-quotient}. Minimality gives
\[
 \frac{\det L_0}{T^{\rank L_0}}
 \le
 \frac{\det\widetilde M}{T^{\rank L_0+\rank M}}
 =
 \frac{\det L_0\,\det M}{T^{\rank L_0+\rank M}},
\]
which proves the quotient bound. Under the additional hypothesis,
writing \(k=\rank L_0\), the two bounds on \(\det L_0\) give
\(e^{-\Delta}T_0^k\le T^k\). Taking logarithms proves the rank
bound.
\end{proof}

\subsection{Robust subsets and small determinants}

Let \(\mathcal A\) be a nonempty finite set with affine span \(F\).
For \(\rho>0\), call \(\mathcal A\) \emph{\(\rho\)-robust} if every
affine subspace \(E\subset F\) of codimension \(t\) satisfies
\[
 |\mathcal A\cap E|\le 2^{-\rho t}|\mathcal A|.
\]
Robustness is preserved by injective affine maps, since their
restrictions to the affine span preserve codimension.

\begin{lemma}[Robust cube extraction]\label{lem:robust-extraction}
Let \(\varnothing\ne\mathcal A\subset\bits^m\), let \(\rho>0\),
and put \(d=\dim\aff\mathcal A\). There is an affine restriction
\(\mathcal A_0=\mathcal A\cap F\), with
\(F=\aff\mathcal A_0\) of dimension \(r\), such that
\(\mathcal A_0\) is \(\rho\)-robust and
\[
 \log_2|\mathcal A_0|
 \ge \log_2|\mathcal A|-\rho(d-r)
 \ge \log_2|\mathcal A|-\rho m,
 \qquad
 \log_2|\mathcal A_0|\le r.
\]
For any fixed \(a_0\in\mathcal A_0\), the number of ordered
\(r\)-tuples \((a_1,\ldots,a_r)\in\mathcal A_0^r\) for which
\(a_1-a_0,\ldots,a_r-a_0\) are independent is at least
\[
 |\mathcal A_0|^r\prod_{t=1}^r(1-2^{-\rho t})
 \ge \bigl((1-2^{-\rho})|\mathcal A_0|\bigr)^r.
\]
When \(r=0\), \(\mathcal A_0\) is a singleton and the number of
empty ordered tuples is one.
\end{lemma}

\begin{proof}
Among all nonempty affine restrictions \(\mathcal Q\) of
\(\mathcal A\), choose one maximizing
\[
 |\mathcal Q|\,2^{-\rho\dim\aff\mathcal Q}.
\]
There are finitely many distinct restrictions. Replacing the
restricting affine space by \(\aff\mathcal Q\) preserves the
restriction, so write the maximizer as
\(\mathcal A_0=\mathcal A\cap F\), \(F=\aff\mathcal A_0\).
Comparison with \(\mathcal A\) gives the lower bounds for its size.

Let \(E\subset F\) have codimension \(t\). If
\(\mathcal A_0\cap E\ne\varnothing\), its affine span has dimension
at most \(r-t\), and the intersection is itself an affine
restriction of \(\mathcal A\). Maximality therefore gives
\[
 |\mathcal A_0\cap E|\,2^{-\rho\dim\aff(\mathcal A_0\cap E)}
 \le |\mathcal A_0|\,2^{-\rho r},
\]
which is the required robustness bound. Empty intersections satisfy
the bound as well. Choose \(r\) coordinate functionals independent
on the direction space of \(F\). Projection to those coordinates is
injective on \(F\), and maps \(\mathcal A_0\) into \(\bits^r\).
Thus \(|\mathcal A_0|\le2^r\).

Finally, after \(i\) independent differences have been chosen, their
affine span with \(a_0\) has codimension \(r-i\) in \(F\).
Robustness leaves at least
\((1-2^{-\rho(r-i)})|\mathcal A_0|\) choices for the next point.
Multiplication for \(i=0,\ldots,r-1\) gives the claimed count. The
empty product proves the rank-zero assertion.
\end{proof}

The next lemma bounds the determinant of the entire difference
lattice. Its affine-image form will also apply to graph lattices:
the only metric input is a bound on image differences.

\begin{lemma}[Determinant of a robust image lattice]\label{lem:robust-height}
Fix \(D>0\) and \(C_0>0\). For all sufficiently large \(N\), put
\(\rho=\ell^{-D}\). Let \(m\le N\), let
\(\varnothing\ne\mathcal A\subset\bits^m\) be \(\rho\)-robust on
its affine span \(F\) of dimension \(r\), and let
\(\Phi:F\longrightarrow\R^d\) be an injective rational affine map,
meaning that it maps rational points of \(F\) to rational points of
\(\R^d\). Suppose that
\[
 \|\Phi(a)-\Phi(b)\|\le C_0\sqrt N
       \qquad(a,b\in\mathcal A).
\]
For any \(a_0\in\mathcal A\), the group
\[
 L_\Phi=\sum_{a\in\mathcal A}\Z\bigl(\Phi(a)-\Phi(a_0)\bigr)
\]
is a rank-\(r\) lattice, independent of the choice of \(a_0\), and
\[
 \det L_\Phi\le \exp\bigl(O_{D,C_0}(N\ell)\bigr).
\]
In particular, with \(W=\Span(\mathcal A-\mathcal A)\),
\[
 H(W)\le \exp\bigl(O_D(N\ell)\bigr).
\]
If \(r=0\), both lattices in question have rank zero and determinant
one.
\end{lemma}

\begin{proof}
The finitely many image differences are rational vectors, so they have
a common denominator. Their integer span is therefore contained in a
scaled copy of \(\Z^d\) and is discrete. The linear part of \(\Phi\) is
injective on the direction space of \(F\), so this lattice has rank
\(r\). Changing \(a_0\)
does not change the generated group. The assertions are immediate
if \(r=0\); suppose \(r\ge1\).

Set \(T=(\log N)^4\), and choose \(L_0\subset L_\Phi\) by
Lemma~\ref{lem:det-minimizer}. Write
\(\Gamma=L_\Phi/L_0\) for its orthogonal quotient and
\(j=\rank\Gamma\). The quotient has
\(\det M\ge T^{\rank M}\) for every sublattice \(M\subset\Gamma\),
while
\(\det L_0\le T^{r-j}\). Fix \(a_0\), and project the vectors
\(\Phi(a)-\Phi(a_0)\) into \(\Gamma\). They have norm at most
\(C_0\sqrt N\). If \(j\ge1\), Lemma~\ref{lem:gaussian-count} at
scale \(T/\tau_j\) bounds the number \(M\) of distinct projections by
\[
 M\le\frac32
     \exp\!\left(\frac{\pi C_0^2N\tau_j^2}{T^2}\right)
 \le 2^{O_{C_0}(N/(\log N)^6)}.
\]
Here \(j\le r\le N\), so \(\tau_j=O(\log N)\). The same final
bound holds for \(j=0\), when \(M=1\).

The preimage in \(F\) of any one projection fiber is an affine
subspace of codimension \(j\). A largest fiber contains at least
\(|\mathcal A|/M\) points, whereas robustness bounds its size by
\(2^{-\rho j}|\mathcal A|\). It follows that
\[
 j\le\frac{\log_2 M}{\rho}
   =O_{C_0}\!\left(\frac{N\ell^D}{(\log N)^6}\right).
\]
The projected differences generate \(\Gamma\). Choose \(j\) of
them that are independent. Their integer span is a full-rank
sublattice of \(\Gamma\), so the index formula and Hadamard's
inequality give
\[
 \det\Gamma\le (C_0\sqrt N)^j.
\]
For \(j=0\) this is the equality \(1=1\). Factoring the determinant
along \(L_0\) now yields
\[
 \det L_\Phi
 \le T^{r-j}(C_0\sqrt N)^j,
 \qquad
 \log\det L_\Phi
 \le 4N\ell+
     O_{C_0}\!\left(\frac{N\ell^D}{(\log N)^5}\right).
\]
For fixed \(D\), the error is \(o(N\ell)\), which proves the stated
determinant bound.

For the identity map the difference norms are at most
\(\sqrt m\le\sqrt N\), and \(L_\Phi\) is a full-rank sublattice
of \(W\cap\Z^m\). Saturating therefore decreases its determinant,
giving \(H(W)\le\det L_\Phi\) and the asserted height bound.
\end{proof}
